\section{The Sample Covariance Problem}

The sample covariance matrix contains \(N(N+1)/2\) distinct variances and covariances. Increasing the number of assets therefore raises the number of parameters at an approximate rate of \(N^2\), without necessarily increasing the available historical observations. Because each estimate contains sampling error, these errors can accumulate as the matrix expands \parencite[378]{ledoit2004well}.

In mathematical terms, the sample covariance matrix, \(S\), is an unbiased estimator of the true covariance matrix, \(\Sigma\) \parencite[368]{ledoit2004well}. Unbiasedness, however, does not imply precision. When the number of historical observations, \(T\), is much larger than the number of assets, \(N\), the sample covariance matrix, \(S\), generally performs adequately. Conversely, when the number of assets is large relative to the number of observations, the sample covariance matrix,
\[
S=\frac{1}{T-1}\sum_{t=1}^{T}(r_t-\bar r)(r_t-\bar r)^{\prime},
\]
performs poorly as an estimator of the true covariance matrix \parencite[366]{ledoit2004well}. After the returns have been demeaned and centred around zero, the sample covariance matrix has a rank of at most \(T-1\). Therefore, when \(N\geq T\), its rank is smaller than its dimension, making the matrix singular and consequently non-invertible \parencite[604]{LEDOIT2003603}. A third possible outcome occurs when \(N<T\), but \(N\) remains close to \(T\), such that the ratio \(N/T\) is large. The sample covariance matrix may then be invertible yet ill-conditioned, meaning that small changes in the estimated covariances can produce large changes in its inverse \parencite[366]{ledoit2004well}.

The accumulation of estimation errors across the matrix also distorts its eigenvalues. The eigenvalues represent the estimated variances of particular combinations of assets. In large-dimensional settings, the largest sample eigenvalues tend to be overestimated, while the smallest tend to be underestimated and pushed closer to zero \parencite[372]{ledoit2004well}. This can result in a drastic amplification of estimation errors when the matrix is inverted, indicating instability in the estimated covariance matrix \parencite[614]{LEDOIT2003603}.

Investors use the inverse covariance matrix to solve the global minimum-variance problem and determine the optimal portfolio weights \parencite[614]{LEDOIT2003603}:
\[
\min_w w^{\prime}S w
\quad\text{subject to}\quad
\mathbf{1}^{\prime}w=1.
\]
Its solution is
\[
\widehat{w}_{\mathrm{GMV}}
=
\frac{S^{-1}\mathbf{1}}
{\mathbf{1}^{\prime}S^{-1}\mathbf{1}}.
\]
Because inversion replaces each eigenvalue with its reciprocal, an apparently low-risk combination associated with a very small eigenvalue can exert disproportionate influence over the weights. The optimiser cannot distinguish the true covariance structure from relationships created by sampling noise and may therefore assign large positive and negative weights to apparent diversification opportunities \parencite[111]{ledoit2004honey}. If these relationships reflect noise, they may not persist out-of-sample, resulting in unstable weights and higher realised risk.
